% GENERATED FILE. Edit canonical lessons, then run npm run generate:books.
% canonical-source-hash: fnv1a64:98b55439efc590c7
% canonical-lessons: RU-C153-sportzal, RU-C153-mechet, RU-C153-khram, RU-C153-pamyatnik, RU-C153-fontan

\chapter{Gym, Mosque, Temple, Monument, Fountain}
\label{ch:ru-gym-mosque-temple-monument-fountain}
\hlchaptermodality{153}

\begin{chapteropening}
\textbf{By the end of this chapter:} \emph{I can say a gym, a mosque, a temple, a monument and a fountain.}

Complete the last lesson of chapter 153: I can say a gym, a mosque, a temple, a monument and a fountain.
\end{chapteropening}

\section[\texorpdfstring{\ru{спортзал} (sportzál)}{спортзал (sportzál)}]{\ru{спортзал} (sportzál) — a gym}

\label{lesson:RU-C153-sportzal}



\begin{quote}
Before the new one: say the Russian for outside, then the Russian for behind.
\end{quote}

\subsection*{You'll want to know: \ru{спортзал}}
\textbf{\ru{спортзал}} (\emph{sportzál}) — \textquotedblleft{}a gym\textquotedblright{}.

\begin{cousinweb}
\textbf{\ru{спортзал}} is \emph{\ru{спорт}}, \textquotedblleft{}sport\textquotedblright{}, from English, with \emph{\ru{зал}}, \textquotedblleft{}a hall\textquotedblright{}, from German \emph{Saal} and French \emph{salle}, words of Germanic origin (Vasmer).
\end{cousinweb}

\begin{grammarlens}[title={What you've built}]
One more word for this chapter.
\end{grammarlens}

\subsection*{Your turn}
\begin{itemize}
\raggedright
  \item \emph{Say it:} \emph{sportzál}
  \item \emph{Say it:} \emph{sportzál}, once more
  \item \emph{Say it:} say \emph{sportzál}
  \item \emph{Recall:} say the Russian for outside, then the Russian for behind, then say \emph{sportzál} again
\end{itemize}

\begin{tcolorbox}[breakable,colback=teal!4,colframe=teal!35!black,title={Before you move on}]
What does \ru{спортзал} mean? (\textquotedblleft{}A gym\textquotedblright{}.) Say it once more.
\end{tcolorbox}
\section[\texorpdfstring{\ru{мечеть} (mechét)}{мечеть (mechét)}]{\ru{мечеть} (mechét) — a mosque}

\label{lesson:RU-C153-mechet}



\begin{quote}
Before the new one: say the Russian for behind, then the Russian for a gym.
\end{quote}

\subsection*{You'll want to know: \ru{мечеть}}
\textbf{\ru{мечеть}} (\emph{mechét}) — \textquotedblleft{}a mosque\textquotedblright{}.

\begin{cousinweb}
\textbf{\ru{мечеть}} — recorded from 1584; it comes from Turkish and Arabic \emph{mäsdžid}, \textquotedblleft{}a place of prostration, a mosque\textquotedblright{}, through Tatar \emph{mäčit} (Vasmer).
\end{cousinweb}

\begin{grammarlens}[title={What you've built}]
One more word for this chapter.
\end{grammarlens}

\subsection*{Your turn}
\begin{itemize}
\raggedright
  \item \emph{Say it:} \emph{mechét}
  \item \emph{Say it:} \emph{mechét}, once more
  \item \emph{Say it:} say \emph{mechét}
  \item \emph{Recall:} say the Russian for behind, then the Russian for a gym, then say \emph{mechét} again
\end{itemize}

\begin{tcolorbox}[breakable,colback=teal!4,colframe=teal!35!black,title={Before you move on}]
What does \ru{мечеть} mean? (\textquotedblleft{}A mosque\textquotedblright{}.) Say it once more.
\end{tcolorbox}
\section[\texorpdfstring{\ru{храм} (khram)}{храм (khram)}]{\ru{храм} (khram) — a temple, a church}

\label{lesson:RU-C153-khram}



\begin{quote}
Before the new one: say the Russian for a gym, then the Russian for a mosque.
\end{quote}

\subsection*{You'll want to know: \ru{храм}}
\textbf{\ru{храм}} (\emph{khram}) — \textquotedblleft{}a temple, a church\textquotedblright{}.

\begin{cousinweb}
\textbf{\ru{храм}} — taken from Church Slavonic, says Vasmer, beside native \emph{\ru{хоромы}}, \textquotedblleft{}a large wooden house\textquotedblright{}: \emph{-\ru{ра}-} against \emph{-\ru{оро}-}. A temple is the house of God.
\end{cousinweb}

\begin{grammarlens}[title={What you've built}]
One more word for this chapter.
\end{grammarlens}

\subsection*{Your turn}
\begin{itemize}
\raggedright
  \item \emph{Say it:} \emph{khram}
  \item \emph{Say it:} \emph{khram}, once more
  \item \emph{Say it:} say \emph{khram}
  \item \emph{Recall:} say the Russian for a gym, then the Russian for a mosque, then say \emph{khram} again
\end{itemize}

\begin{tcolorbox}[breakable,colback=teal!4,colframe=teal!35!black,title={Before you move on}]
What does \ru{храм} mean? (\textquotedblleft{}A temple, a church\textquotedblright{}.) Say it once more.
\end{tcolorbox}
\section[\texorpdfstring{\ru{памятник} (pámyatnik)}{памятник (pámyatnik)}]{\ru{памятник} (pámyatnik) — a monument}

\label{lesson:RU-C153-pamyatnik}



\begin{quote}
Before the new one: say the Russian for a mosque, then the Russian for a temple.
\end{quote}

\subsection*{You'll want to know: \ru{памятник}}
\textbf{\ru{памятник}} (\emph{pámyatnik}) — \textquotedblleft{}a monument\textquotedblright{}.

\begin{cousinweb}
\textbf{\ru{памятник}} is \emph{\ru{память}}, \textquotedblleft{}memory\textquotedblright{}, with \emph{-\ru{ник}}: a thing that keeps memory. Vasmer relates \emph{\ru{память}} to Lithuanian \emph{atmintìs}, \textquotedblleft{}memory\textquotedblright{}, Sanskrit \emph{matí-}, \textquotedblleft{}thought\textquotedblright{}, and Latin \emph{mēns}, \textquotedblleft{}mind\textquotedblright{}.
\end{cousinweb}

\begin{grammarlens}[title={What you've built}]
One more word for this chapter.
\end{grammarlens}

\subsection*{Your turn}
\begin{itemize}
\raggedright
  \item \emph{Say it:} \emph{pámyatnik}
  \item \emph{Say it:} \emph{pámyatnik}, once more
  \item \emph{Say it:} say \emph{pámyatnik}
  \item \emph{Recall:} say the Russian for a mosque, then the Russian for a temple, then say \emph{pámyatnik} again
\end{itemize}

\begin{tcolorbox}[breakable,colback=teal!4,colframe=teal!35!black,title={Before you move on}]
What does \ru{памятник} mean? (\textquotedblleft{}A monument\textquotedblright{}.) Say it once more.
\end{tcolorbox}
\section[\texorpdfstring{\ru{фонтан} (fontán)}{фонтан (fontán)}]{\ru{фонтан} (fontán) — a fountain}

\label{lesson:RU-C153-fontan}



\begin{quote}
Before the new one: say the Russian for a temple, then the Russian for a monument.
\end{quote}

\subsection*{You'll want to know: \ru{фонтан}}
\textbf{\ru{фонтан}} (\emph{fontán}) — \textquotedblleft{}a fountain\textquotedblright{}.

\begin{cousinweb}
\textbf{\ru{фонтан}} — from Peter the Great's time; it came through Polish \emph{fontana} or straight from Italian \emph{fontana}, from Latin \emph{fōns}, \emph{fontis}, \textquotedblleft{}a spring\textquotedblright{} (Vasmer).
\end{cousinweb}

\begin{grammarlens}[title={What you've built}]
One more word for this chapter.
\end{grammarlens}

\subsection*{Your turn}
\begin{itemize}
\raggedright
  \item \emph{Say it:} \emph{fontán}
  \item \emph{Say it:} \emph{fontán}, once more
  \item \emph{Say it:} say \emph{fontán}
  \item \emph{Recall:} say the Russian for a temple, then the Russian for a monument, then say \emph{fontán} again
\end{itemize}

\begin{tcolorbox}[breakable,colback=teal!4,colframe=teal!35!black,title={Before you move on}]
What does \ru{фонтан} mean? (\textquotedblleft{}A fountain\textquotedblright{}.) Say it once more.
\end{tcolorbox}
